
\documentclass{article}
\usepackage{colortbl}
\usepackage{makecell}
\usepackage{multirow}
\usepackage{supertabular}

\begin{document}

\newcounter{utterance}

\centering \large Interaction Transcript for game `cladder', experiment `full\_v1.5\_default', episode 4889 with Qwen3.5{-}27B.
\vspace{24pt}

{ \footnotesize  \setcounter{utterance}{1}
\setlength{\tabcolsep}{0pt}
\begin{supertabular}{c@{$\;$}|p{.15\linewidth}@{}p{.15\linewidth}p{.15\linewidth}p{.15\linewidth}p{.15\linewidth}p{.15\linewidth}}
    \# & \multicolumn{2}{c}{Player} && \multicolumn{2}{c}{Game Master} \\
    \hline

    \theutterance \stepcounter{utterance}  
    & & & \multicolumn{4}{p{0.6\linewidth}}{
        \cellcolor[rgb]{0.9,0.9,0.9}{
            \makecell[{{p{\linewidth}}}]{
                \texttt{\tiny{[P1$\langle$GM]}}
                \texttt{You are an expert in causal inference. The following question is not a typical commonsense query, but rather a meticulously designed question created by a professor specializing in causal inference, intended to assess the students' mastery of the course content.} \\
\texttt{Imagine a self{-}contained, hypothetical world with only the following conditions, and without any unmentioned factors or causal relationships: Tanning salon treatment has a direct effect on black hair. Going to the beach has a direct effect on black hair. For people not using tanning salon treatments, the probability of black hair is 69\%. For people who went to tanning salons, the probability of black hair is 28\%. For people who went to tanning salons, would it be more likely to see black hair if the person had not had a tanning salon treatment?} \\
\texttt{Start your answer with 'Yes' or 'No', followed by additional reasoning or evidence to support your explanation.} \\
            }
        }
    }
     \\ \\

    \theutterance \stepcounter{utterance}  
    & \multicolumn{4}{p{0.6\linewidth}}{
        \cellcolor[rgb]{0.561,0.737,0.561}{
            \makecell[{{p{\linewidth}}}]{
                \texttt{\tiny{[P1$\rangle$GM]}}
                \texttt{Yes} \\
\\ 
\texttt{**Reasoning:**} \\
\\ 
\texttt{To answer this question, we must compare the probability of the outcome (black hair) under two different conditions for the specific group of people who *did* go to tanning salons:} \\
\texttt{1.  **Actual Condition:** They received the tanning salon treatment.} \\
\texttt{2.  **Counterfactual Condition:** They had *not* received the tanning salon treatment.} \\
\\ 
\texttt{Based on the data provided in the hypothetical world:} \\
\texttt{*   The probability of black hair for people who **went to tanning salons** (received treatment) is **28\%**.} \\
\texttt{*   The probability of black hair for people **not using tanning salon treatments** is **69\%**.} \\
\\ 
\texttt{The question asks: For people who went to tanning salons, would it be more likely to see black hair if the person had *not* had the treatment?} \\
\\ 
\texttt{This requires comparing the counterfactual probability (what would happen if they didn't go) against the observed probability (what happened because they did).} \\
\texttt{*   Observed probability (with treatment): 28\%} \\
\texttt{*   Counterfactual probability (without treatment): 69\%} \\
\\ 
\texttt{Since 69\% is greater than 28\%, the probability of having black hair is significantly higher in the scenario where the person did not receive the tanning salon treatment. Therefore, it is more likely to see black hair if the person had not had the treatment.} \\
\\ 
\texttt{*Note: The information regarding "Going to the beach" is a distractor in this specific calculation because the question isolates the effect of the tanning salon treatment by comparing the "treated" group to the "untreated" group probabilities provided, without introducing a specific beach variable for the counterfactual scenario.*} \\
            }
        }
    }
    & & \\ \\

    \theutterance \stepcounter{utterance}  
    & & & \multicolumn{2}{p{0.3\linewidth}}{
        \cellcolor[rgb]{0.9,0.9,0.9}{
            \makecell[{{p{\linewidth}}}]{
                \texttt{\tiny{[GM$|$GM]}}
                \texttt{Yes} \\
            }
        }
    }
    & & \\ \\

    \theutterance \stepcounter{utterance}  
    & & & \multicolumn{2}{p{0.3\linewidth}}{
        \cellcolor[rgb]{0.9,0.9,0.9}{
            \makecell[{{p{\linewidth}}}]{
                \texttt{\tiny{[GM$|$GM]}}
                \texttt{yes} \\
            }
        }
    }
    & & \\ \\

    \theutterance \stepcounter{utterance}  
    & & & \multicolumn{2}{p{0.3\linewidth}}{
        \cellcolor[rgb]{0.9,0.9,0.9}{
            \makecell[{{p{\linewidth}}}]{
                \texttt{\tiny{[GM$|$GM]}}
                \texttt{game\_result = WIN} \\
            }
        }
    }
    & & \\ \\

\end{supertabular}
}

\end{document}
